\section{Related work}
\textbf{Body-part classification.} Pham et al. (2021) sorted 16,093 images into five coarse groups for DICOM routing (recall 0.982) and pointed out that routing needs near-perfect accuracy. Langer et al. (2024, preprint) classified 14 regions on 48,434 hospital radiographs (96.6 \% for a single model, 97.7 \% for an ensemble). When experts reviewed the test errors, 35 \% turned out to be wrong labels and 11 \% were out-of-domain images, so label noise accounts for part of the remaining error.

Mitsuyama et al. (2026) trained a seven-category model (Head, Neck, Chest, Incomplete Chest, Abdomen, Pelvis, Extremities) on 429,341 radiographs from two institutions and MURA, and tested it on a second institution and on MURA. Every MURA image was labelled "Extremities", and micro-accuracy on MURA was 98.6 \% (99 \% CI 98.4 to 98.7 \%). Radiologist review changed the DICOM-derived body-part label of 3.9 \% of the radiographs from one institution and 12.8 \% from the other. Their MURA figure cannot be compared with ours: theirs is a coarse seven-way task in which all MURA images share one class, whereas we separate the seven MURA study types.

On MURA, Shah et al. (2026) classified the seven upper-extremity study types with an attention-fused Xception-Swin model and reached an accuracy of 0.9643, a macro-F1 of 0.9574 and an AUC of 0.9963. They used the public validation set (3,197 images) as the test partition, a fixed seed (42), a patient-wise data organisation and image-level evaluation, and they report Wilson confidence intervals for the abnormality-detection accuracies; we did not find intervals or repeated seeds for the body-part result. All of these studies differ from ours in data, class set and metric.

\textbf{Hybrid CNN-Transformer models} with gated or attention-based fusion are common. We did not find a report that inspects the learned gate, although absence from our search does not prove absence from the literature. \textbf{Hierarchical losses} have also been tried: Wu, Tygert and LeCun (2019) found that minimising a hierarchical loss directly is difficult, and that minimising cross-entropy often lowers it almost as well, which makes it more useful as a metric than as a training objective.

\textbf{Evaluation quality.} Leakage is widespread in machine-learning-based science (Kapoor and Narayanan, 2023) and in COVID-19 imaging (Roberts et al., 2021). Label quality limits medical-imaging benchmarks (Oakden-Rayner, 2020), and reporting guidance exists in the CLAIM 2024 update (Tejani et al., 2024). Oakden-Rayner (2020) audited MURA labels and found them more accurate than those of ChestX-ray14, although the normal/abnormal labels were unreliable for cases with degenerative joint disease. Confident-learning methods find label errors from out-of-fold probabilities (Northcutt et al., 2021).


